\chapter{Circle and Trigonometry}\label{ch:rational-circle-trigonometry}

Circle geometry supplies the first special function without assuming
standard real trigonometry.  Exact rational coordinates give squared
identities; length and angle computations are added where needed.  The
sine convention throughout is
\[
 S(x)=\sin(\pi x),\qquad C(x)=\cos(\pi x).
\]
These names designate the constructions below, not functions imported with
an existing theory of differentiation.

\section{Finite rational geometry}
\label{thm:rational-circle-pythagorean}
The rational chart
\[
 P(t)=\left(\frac{1-t^2}{1+t^2},\frac{2t}{1+t^2}\right)
\]
satisfies $P_1(t)^2+P_2(t)^2=1$ by expansion.  For $0\le t\le1$ it runs
through the first quadrant, from $(1,0)$ to $(0,1)$.  Conversely a point of
the unit circle in this quadrant has parameter $t=y/(1+x)$.

Rational determinants give triangle and polygon areas.  The broken-line
inequality of Chapter~\ref{ch:foundations} compares polygonal lengths.
Reflections and the rational rotation matrix
\[
 \begin{pmatrix}c&-s\\s&c\end{pmatrix},\qquad c^2+s^2=1,
\]
preserve squared lengths and determinants by finite algebra.  These are all
the geometric operations used in this chapter.

\section{Archimedes and rational circles}
\label{thm:geometric-pi-validity}
Join finitely many points $P(t_i)$, where
$0=t_0<\cdots<t_N=t\le1$, to the origin.  The inscribed triangle for
$u=t_i,v=t_{i+1}$ has area
\[
 I(u,v)=\frac{(v-u)(1+uv)}{(1+u^2)(1+v^2)}.
\]
The two endpoint tangents meet outside the circle; the quadrilateral formed
with the origin has area
\[
 O(u,v)=\frac{v-u}{1+uv}.
\]
Both formulas follow by determinants and the tangent equations
$P(u)\mathbin{\cdot}z=P(v)\mathbin{\cdot}z=1$.

Inserting a point increases the total inner area and decreases the total
outer area: the changes are areas of finite triangles.  Moreover
\[
 \frac{v-u}{1+v^2}\le I(u,v)\le O(u,v)
       \le\frac{v-u}{1+u^2}.
\]
Consequently the gap between the total outer and inner areas is at most
\[
 \mesh(P)\left(1-\frac1{1+t^2}\right)\le\mesh(P).
\]
Nested rational subdivisions therefore define a sector-area computation
$A(t)$.  Define
\[
 \pi=4A(1).
\]
The square and a finer inscribed polygon give $2<\pi<4$.  The strict
bounds are witnessed by finite area computations.  No circle integral or
series formula has entered the definition.

\begin{figure}[htbp]\centering
\rationalCircleSubdivisionAnimation
\caption{Rational subdivisions give inner and outer bounds for the circle.}
\label{fig:rational-parameter-line}
\end{figure}

Area and circumference are initially independent computations.  Here is
their comparison.  For an inscribed chord of length $c$, its distance from
the origin is $d=\sqrt{1-c^2/4}$, and twice the corresponding triangle area
is $cd$.  Thus
\[
 0\le c-cd\le c^3/4.
\]
The total inscribed length is bounded by the outer tangent length.  For a
unit-circle tangential polygon, twice its area equals its perimeter,
triangle by triangle.  As the largest chord shrinks, the displayed bound
identifies circumference with twice disk area.  Hence the circumference is
$2\pi$.  The same proof works for sectors: arc length is twice sector area.
Rotation invariance and additivity follow by common refinements of the
finite polygons.  They remain true for computed rotation matrices by
coordinate approximation and the determinant error bounds.

\section{Angles and rotations}
\label{thm:finite-quarter-turn}
A full turn has normalized parameter length $2$; a quarter turn has length
$1/2$.  The point whose sector area is $\pi x/2$ is denoted $(C(x),S(x))$.
The inverse construction below gives it for every rational
$0\le x\le1/2$.  Reflection and periodicity then give every rational input.

Equal dyadic subdivisions offer another computation.  If $(c,s)$ is the
parent point in the first quadrant, the half-angle point is
\[
 \left(\sqrt{(1+c)/2},\sqrt{(1-c)/2}\right).
\]
Squaring and applying the double-angle rotation identity proves this is its
positive square root as a rotation.  Repeated halving and finite rotation
products compute all dyadic angles.  Additivity of sector area shows that
halving the rotation halves the angle parameter, so these are the same
points as the sector-area construction.  The comparison is geometric,
not an appeal to a completed trigonometric function.

Rotation multiplication now gives
\begin{align*}
 S(x+y)&=S(x)C(y)+C(x)S(y),\\
 C(x+y)&=C(x)C(y)-S(x)S(y),\\
 C(x)^2+S(x)^2&=1.
\end{align*}
The endpoint values are
$S(0)=0$, $C(0)=1$, $S(1/2)=1$, $C(1/2)=0$.
Symmetries give $S(-x)=-S(x)$, $C(-x)=C(x)$, and period $2$.

Chord length is at most arc length.  Thus, for rational arguments,
\[
 |S(x)-S(y)|\le4|x-y|,\qquad
 |C(x)-C(y)|\le4|x-y|.
\]
This gives a rational continuity certificate.  In particular dyadic
approximation evaluates every rational angle, and all these representations
agree by the sector-area comparison.

\section{Arctangent as the angle chart}
For $0\le u<v\le1$, the area bounds already proved give
\[
 \frac{v-u}{1+v^2}\le A(v)-A(u)
       \le\frac{v-u}{1+u^2}.
\]
In particular
\[
 \frac{v-u}{2}\le A(v)-A(u)\le v-u.
\]
The chart is interval-continuous and has a linear inverse-separation bound.
For rational $x\in[0,1/2]$, the target $\pi x/2=2xA(1)$ is certified
between its endpoint values.  Trisection therefore computes the unique
parameter $t$ with
\[
 A(t)=\pi x/2.
\]
Evaluate $P(t)$ by the computed-input construction to obtain $C(x),S(x)$.
At the endpoints use $t=0,1$ directly.

The notation $A(t)=\arctan t$ records the usual angle chart: the geometric
angle of $P(t)$ is $2A(t)$.  On a partition of $[0,t]$, the rectangle
sums for $1/(1+u^2)$ bracket the same inner and outer areas.  Their gap is
at most the mesh.  Thus these rectangles are already an equivalent
computation of $A(t)$; Chapter~\ref{ch:integrals} will recognize them as
an integral.  This order of construction avoids defining the sine through
an FTC not yet proved.

\section{Small angles before differentiation}
\label{sec:geometric-small-angle}
For $0<h\le1/4$, compare the inscribed triangle, the sector, and the
triangle cut out by its tangent.  Using arc length $\pi h$ gives
\[
 S(h)\le\pi h\le S(h)/C(h).
\]
The half-angle identity and the chord bound also give
\[
 0\le1-C(h)=2S(h/2)^2\le8h^2.
\]
Multiplying the first comparison by the positive $C(h)$ and using
$\pi<4$ yields
\[
 0\le\pi h-S(h)\le\pi h(1-C(h))\le32h^3.
\]
Odd and even symmetry extend the estimates to negative $h$:
\[
 |S(h)-\pi h|\le32|h|^3,\qquad |1-C(h)|\le8h^2.
\]
These are finite geometric error bounds.  They will prove both the
termination of derivative computations and their identification with
$\pi C$ and $-\pi S$.

\section{Input to integration}
The first area to be computed in two ways is
\[
 \int_0^{1/2}S(x)\,dx.
\]
Equal dyadic samples use nested square roots.  Unequal samples use the
rational parameter $t$ and the chart $P(t)$.  Both are legitimate finite
computations; their integral equivalence follows from interval continuity
and refinement.  The derivative computation and the FTC will identify the
answer as $1/\pi$.

Only geometry that certifies these infinite processes has been used.
Natural charts may be rational coordinates, dyadic angles, or computed
square roots.  The mathematical function is tracked through proved
comparisons between these descriptions, not through a preferred closed
formula.
